\chapter{Conclusion}
\label{chap:conclusion}

This thesis investigates a staged optimization process combining serving configuration, speculative decoding, quantization and targeted GPU execution changes. The intended contribution is evidence about which combinations improve end-to-end performance under defined quality and completion constraints.

The completed work establishes bilingual workloads, a benchmark procedure and 36 Qwen3-8B BF16 baseline measurements on one H100. The analysis identifies variable response lengths, an anomalous concurrency-one run and substantial cancellations in one high-load condition. These limitations require validation before final comparisons.

The remaining work is to repeat the baseline, evaluate the prepared DFlash integration with matched controls, and test supported precision and engine configurations. Profiling, component-removal experiments and held-out quality evaluation will support the final conclusions.

\pending{Final conclusion: insert measured answers to RQ1--RQ3, supported performance gains, quality and completion constraints, and remaining limitations.}
